%----------
%	SOLUCIÓN TECNOLÓGICA BASADA EN IA (intro de 3.2)
%----------

La tarea descrita en la Sección~\ref{sec:experts} pide clasificar cada pieza
periodística según cinco variables de sexismo lingüístico. Dos rasgos de esa
tarea determinan la aproximación tecnológica. El primero: no hay un conjunto de
entrenamiento etiquetado grande, porque la anotación experta es cara y cubre
1.313 piezas. El segundo: codificar no consiste en buscar palabras clave, sino
en aplicar una metodología interpretativa (definiciones, reglas de decisión,
fronteras entre variables) tomada del libro de códigos del proyecto.

Estos dos rasgos desaconsejan el aprendizaje supervisado clásico, que necesitaría
miles de ejemplos etiquetados. En su lugar se recurre a modelos generativos de
lenguaje (LLM) en régimen \textit{zero-shot}: el modelo no aprende la tarea a
partir de ejemplos, sino que recibe la metodología en el propio \textit{prompt} y
la aplica sobre cada texto. El libro de códigos hace de supervisión en tiempo de
inferencia, sin reentrenar nada.

La solución es una tubería común, independiente del proveedor. Por cada pieza y
variable, el sistema construye una entrada con la metodología y el texto, la
envía a un modelo y recupera una salida estructurada con tres campos: un código,
una explicación razonada y las evidencias literales que la sustentan. Esa misma
tubería se enruta, según el identificador del modelo, tanto a modelos locales
como a las API de OpenAI, Google y Anthropic (Sección~\ref{sec:llms_models_iris}).
Así, las diferencias que se observen entre modelos no se confunden con
diferencias de implementación.

\begin{figure}[H]
\centering
\begin{tikzpicture}[
  node distance=8mm and 12mm,
  caja/.style={draw, rounded corners=3pt, align=center, minimum height=12mm,
               inner sep=5pt, font=\small},
  io/.style={caja, fill=black!5},
  proc/.style={caja, fill=black!2},
  arr/.style={-{Latex[length=2.4mm]}, thick, gray}
]
  \node[io] (art) {Artículo};
  \node[proc, right=of art] (prompt) {Metodología\\ + texto};
  \node[proc, right=of prompt] (modelo) {Modelo\\ (LLM)};
  \node[io, right=of modelo, align=left] (out)
        {\texttt{codigo}\\ \texttt{explicacion}\\ \texttt{evidencias}};
  \draw[arr] (art) -- (prompt);
  \draw[arr] (prompt) -- (modelo);
  \draw[arr] (modelo) -- (out);
\end{tikzpicture}
\caption{Tubería de clasificación. Por cada pieza y variable, la metodología del
libro de códigos y el texto forman la entrada del modelo, que devuelve una salida
estructurada y auditable.}
\label{fig:iris-pipeline}
\end{figure}

La salida estructurada responde al marco \textit{human-in-the-loop} del proyecto.
Pedir explicación y evidencias no es un capricho de formato: convierte cada
predicción en algo que la analista puede auditar, aceptar o rechazar. El sistema
asiste; no decide en su lugar.

Las secciones siguientes describen los modelos generativos considerados
(\ref{sec:llms_models_iris} y \ref{sec:genai_models_iris}) y el paradigma
\textit{zero-shot} con libro de códigos (\ref{sec:zeroshot_iris}). Las dos
arquitecturas de inferencia que se comparan, inyectar la metodología en el
\textit{prompt} frente a empaquetarla como \textit{Agent Skills}, se detallan en
el Capítulo~\ref{implementation}.
